% Problem record op_452fc7d8cc44b728. This TeX file is derived from
% database/problems_json/op_452fc7d8cc44b728.json by scripts/sync-tex.mjs; edit the JSON record
% and rerun the script rather than editing this file.
\section{Universal simulation with two PR boxes}
\label{sec:op_452fc7d8cc44b728}

\paragraph{Problem.}
Can shared randomness and at most two Popescu--Rohrlich boxes exactly simulate
every pair of local projective measurements on every two-qubit state, without
communication?  A Popescu--Rohrlich box is the binary nonsignalling behavior
\begin{equation}
  P_{\mathrm{PR}}(u,v\mid s,t)
  =\begin{cases}
     \tfrac12,&u\oplus v=st,\\
     0,&u\oplus v\neq st,
   \end{cases}
  \qquad s,t,u,v\in\{0,1\}.
  \label{eq:p28-pr-box}
\end{equation}
The parties may use arbitrary local, possibly adaptive wirings of two
independent copies of Eq.~\eqref{eq:p28-pr-box}.  By convexity and Schmidt
decomposition, it suffices to solve the simulation for all pure states
\begin{equation}
  \lvert\psi_\theta\rangle
  =\cos\theta\,\lvert00\rangle+\sin\theta\,\lvert11\rangle,
  \qquad 0<\theta\leq\frac{\pi}{4},
  \label{eq:p28-schmidt-state}
\end{equation}
and all local projective measurements on the state in
Eq.~\eqref{eq:p28-schmidt-state}.  A complementary finite-scenario objective
is to characterize the convex set generated by two-box wirings, for example
through its facet inequalities.

\subsection*{Status}

Unsolved

\subsection*{Source}

Gisin explicitly asks whether two Popescu--Rohrlich boxes suffice to simulate
all projective-measurement correlations of arbitrary two-qubit pure states
\sourcecite{ref:p28-gisin}{Gis09}.

\subsection*{Progress}

\begin{itemize}
  \item Gisin explicitly posed both the universal two-box simulation problem
  and the finite-scenario inequality problem
  \sourcecite{ref:p28-gisin}{Gis09}.

  \item One copy of Eq.~\eqref{eq:p28-pr-box}, together with shared
  randomness, exactly simulates arbitrary projective measurements on a
  maximally entangled qubit pair \sourcecite{ref:p28-cerf}{CGMP05}.  This
  covers only the endpoint $\theta=\pi/4$ in
  Eq.~\eqref{eq:p28-schmidt-state}.

  \item Finite-setting inequalities valid for all one-PR-box strategies are
  violated by correlations of nonmaximally entangled two-qubit states.
  Hence one PR box is not a universal resource, but these inequalities do not
  decide whether two boxes suffice
  \sourcecite{ref:p28-brunner-njp}{BGS05},
  \sourcecite{ref:p28-brunner-jmp}{BSG06}.

  \item Photonic experiments have violated one-box bounds and thereby
  certified a two-box lower bound for selected measurements.  This does not
  construct a two-box simulation valid for every state and measurement
  \sourcecite{ref:p28-christensen}{CLBGK15}.
\end{itemize}

\subsection*{References}

\begin{enumerate}[label={},leftmargin=4.8em,labelsep=0.5em]
  \footnotesize
  \item[\textup{[Gis09]}]\label{ref:p28-gisin}
  N. Gisin, ``Bell Inequalities: Many Questions, a Few Answers,'' in
  W. C. Myrvold and J. Christian (eds.), \emph{Quantum Reality, Relativistic
  Causality, and Closing the Epistemic Circle}, The Western Ontario Series in
  Philosophy of Science \textbf{73}, 125--138 (Springer, 2009).
  \href{https://doi.org/10.1007/978-1-4020-9107-0_9}{doi:10.1007/978-1-4020-9107-0\_9};
  \href{https://arxiv.org/abs/quant-ph/0702021}{arXiv:quant-ph/0702021}.

  \item[\textup{[CGMP05]}]\label{ref:p28-cerf}
  N. J. Cerf, N. Gisin, S. Massar, and S. Popescu, ``Simulating Maximal
  Quantum Entanglement without Communication,'' \emph{Physical Review
  Letters} \textbf{94}, 220403 (2005).
  \href{https://doi.org/10.1103/PhysRevLett.94.220403}{doi:10.1103/PhysRevLett.94.220403};
  \href{https://arxiv.org/abs/quant-ph/0410027}{arXiv:quant-ph/0410027}.

  \item[\textup{[BGS05]}]\label{ref:p28-brunner-njp}
  N. Brunner, N. Gisin, and V. Scarani, ``Entanglement and Non-locality Are
  Different Resources,'' \emph{New Journal of Physics} \textbf{7}, 88
  (2005). \href{https://doi.org/10.1088/1367-2630/7/1/088}{doi:10.1088/1367-2630/7/1/088};
  \href{https://arxiv.org/abs/quant-ph/0412109}{arXiv:quant-ph/0412109}.

  \item[\textup{[BSG06]}]\label{ref:p28-brunner-jmp}
  N. Brunner, V. Scarani, and N. Gisin, ``Bell-Type Inequalities for
  Non-local Resources,'' \emph{Journal of Mathematical Physics} \textbf{47},
  112101 (2006).
  \href{https://doi.org/10.1063/1.2352857}{doi:10.1063/1.2352857};
  \href{https://arxiv.org/abs/quant-ph/0603094}{arXiv:quant-ph/0603094}.

  \item[\textup{[CLBGK15]}]\label{ref:p28-christensen}
  B. G. Christensen, Y.-C. Liang, N. Brunner, N. Gisin, and P. G. Kwiat,
  ``Exploring the Limits of Quantum Nonlocality with Entangled Photons,''
  \emph{Physical Review X} \textbf{5}, 041052 (2015).
  \href{https://doi.org/10.1103/PhysRevX.5.041052}{doi:10.1103/PhysRevX.5.041052};
  \href{https://arxiv.org/abs/1506.01649}{arXiv:1506.01649}.
\end{enumerate}

\subsection*{Comment}

Neither a universal exact two-box protocol nor a counterexample requiring
more than two boxes is known.  In finite input--output scenarios, a complete
description of the convex set generated by two-box wirings is likewise open.

\subsection*{Field}

Quantum foundations

\subsection*{Topic}

Bell nonlocality; Resource conversion; Qubit systems

\subsection*{ID}

\texttt{op\_452fc7d8cc44b728}
